% ===================================================================
\section{Chemical Data Governance and Cross-Registry Linking}
% ===================================================================

\subsection{The Chemical Data Challenge}

A particularly challenging category of \DPP\ data is chemical composition and hazard information. The \ESPR, \REACH, the Battery Regulation, and product-specific delegated acts all require that \DPPs\ carry data on substances of concern, hazardous classifications, and material composition. However, this data category sits at the intersection of three structural tensions: (1)~the regulatory mandate for transparency, (2)~the commercial sensitivity of proprietary formulations, and (3)~the scientific authority of established chemical registries. A representative from the materials industry captured this tension directly at the \DPPforEU~2026 conference~\cite{dpp4eu2026_day2}, noting that the \DPP\ must address ``\textit{material safety, our chemical safety processes}'' while also acknowledging that there is ``\textit{a lot more that we need from this community in order to meet us in the middle}''. The distributed flat-file model resolves none of these tensions: it either exposes proprietary data or withholds it, and it either duplicates chemical hazard information from authoritative sources or omits it entirely.

\subsection{The Duplication Problem}

Chemical hazard data --- substance identities, hazard classifications, \SVHC\ listings, safety data sheet content --- is already maintained by authoritative public databases. The European Chemicals Agency (ECHA)~\cite{echa} maintains the authoritative registry of substances registered under \REACH, including \SVHC\ candidate list entries, harmonised classifications under \CLP, and authorisation list entries. \PubChem~\cite{pubchem} maintained by the U.S. National Library of Medicine provides open-access chemical substance information including structures, properties, and links to regulatory entries. Duplicating this data within individual \DPPs\ creates three problems. First, \textbf{data drift}: when ECHA updates a hazard classification or adds a substance to the \SVHC\ candidate list, all \DPPs\ that have duplicated the old classification become stale and inconsistent. Second, \textbf{scale}: with potentially millions of \DPPs\ across the European single market, each carrying duplicated chemical data, the storage and update burden becomes prohibitive. Third, \textbf{authority}: a hazard classification stated in a \DPP\ carries less authority than a classification linked to the ECHA database, because the \DPP\ data is self-declared by the \REO\ while the ECHA data is regulatorily authoritative.

\subsection{Cross-Registry Linking via External Entity Normalisation}

The portal-backed architecture resolves the duplication problem through the external entity normalisation stage of the Ingestion Governance Loop (Section~4.2.3). When an \REO\ submits a \DPP\ containing chemical substance references, the portal controller normalises these references to canonical \IRIs\ using external registry lookup services (e.g., ``\CAS\ 7440-44-0'' to \texttt{urn:echa:substance:7440-44-0}). The \DPP\ stores the \textit{link}, not the \textit{data}. When a user accesses the \DPP\ through the Resolution Governance Loop, the portal can resolve the link in real time, retrieving the current authoritative hazard classification from ECHA and presenting it alongside the \REO's composition declaration. This approach ensures that: (1)~the \DPP\ always reflects the most current authoritative hazard data, (2)~the storage burden is minimised (only links, not full data sets), and (3)~the authority of the data is preserved (it comes from ECHA, not from the \REO's self-declaration). The conference discussion on material declaration reinforced this principle, with a participant advocating for ``single point of definition'' (SPOD) --- ``\textit{one repository where we say that's how we identify a material and that are the definitions we serve to tell you what are the hazard statements what is the granularity}''~\cite{dpp4eu2026_day1}. The cross-registry linking approach operationalises the \SPOD\ principle within the portal-backed architecture.

\subsection{The Quality Infrastructure Connection}

The cross-registry linking strategy extends beyond chemical databases to the broader quality infrastructure: standards, measurements, and certificates. As a representative of the German quality infrastructure emphasised at the \DPPforEU~2026 conference~\cite{dpp4eu2026_day2}, the quality infrastructure ``\textit{allows to link the data from the \DPP\ to exactly those standards, measurements and certificates}'', with the further observation that certificate data could be directly implemented in the \DPP\ and that it is important to trace what standards were used to assess the information~\cite{dpp4eu2026_day2}. The portal-backed architecture realises this vision through external entity normalisation: when a \DPP\ references a standard (e.g., ``EN 18215''), a measurement (e.g., a specific test method), or a certificate (e.g., a conformity assessment certificate from an accredited body), the portal controller normalises these references to canonical \IRIs\ linking to the relevant quality infrastructure databases. The conference discussion also highlighted the importance of the quality infrastructure for trust: ``\textit{I think we have many \DPP\ proposals trying to make verifiable data without addressing, in particular, the QI, the quality infrastructure, and I think it's a waste of time because this is what we need to have in order to trust the \DPP\ structures at all}''~\cite{dpp4eu2026_day2}. The three-pillar trust tier matrix (Section~5.5) complements this by ensuring that the organisations providing certified data are themselves verified through \EORI\ / \VIES, \WthreeC\ Verifiable Credentials, and \eIDAS~2.0 identity binding.

\subsection{The \SVHC\ and \CLP\ Alignment}

A specific challenge is the representation of Substances of Very High Concern (SVHC) and harmonised classifications under the Classification, Labelling and Packaging (CLP) regulation in \DPPs. A conference participant from the materials manufacturing perspective noted at \DPPforEU~2026~\cite{dpp4eu2026_day1} that the industry needs consistency: ``\textit{if we talk about \SVHCs, follow the same definition from \CLP\ and the candidate list from ECHA whether you work in textiles, automotive tyres, or furniture}''. The cross-registry linking approach ensures this consistency structurally: rather than each \REO\ defining \SVHC\ and \CLP\ classifications independently in their \DPPs, all \DPPs\ link to the same ECHA-maintained candidate list and \CLP\ harmonised classification entries. When ECHA adds a substance to the \SVHC\ candidate list, all \DPPs\ containing that substance automatically reflect the updated status upon resolution, without requiring any action from the \REOs. This design directly addresses the conference concern about consistency and eliminates the risk of divergent \SVHC\ classifications across product categories.

\subsection{The Industry Dilemma: Proprietary Composition and Partial Disclosure}

A well-documented tension in \DPP\ implementation is the resistance of chemical and materials companies to disclosing proprietary composition data through a system that could expose it to competitors. The conference discussion on confidential data sharing~\cite{dpp4eu2026_day1} exemplifies this concern. The \ABAC\ framework with \ODRL\ branch-level policies (Section~5) provides the architectural answer: proprietary composition branches are protected by \ODRL\ policies that restrict access to 
\\
\texttt{dpp:InternalDataOriginator} and 
\\
\texttt{dpp:MarketSurveillanceAuthority}, 
\\
while public-facing composition summaries (which may omit proprietary details) are accessible to \texttt{dpp:PublicUser}. The cross-registry linking strategy adds a further dimension: for hazard classification data that is already public in ECHA, the \DPP\ can link to ECHA rather than duplicating the classification, allowing the public to see the hazard information without the \REO\ disclosing its proprietary formulation. The \REO\ declares the presence of the substance (and links to its ECHA entry) without disclosing the exact quantity or formulation context. This partial disclosure model --- link to public hazard data, protect proprietary composition --- resolves the Industry dilemma within a single architectural framework. The conference observation that ``\textit{this data really can be shared and if there's a compromise between confidentiality and transparency}''~\cite{dpp4eu2026_day1} is realised through this granular, branch-level approach.

\subsection{Interoperability with Industry Data Platforms}

The cross-registry linking strategy must also address interoperability with industry-specific data platforms that already maintain chemical and materials data. A conference presentation demonstrated interoperability between ChemX and ContainerX at the Hannover Messe, where ``\textit{SECA as a client picked up the data from Covestro that was polymer data}''~\cite{dpp4eu2026_day2}. The portal-backed architecture supports such interoperability through external entity normalisation: when a \DPP\ references a material from an industry platform (e.g., a Covestro polymer grade), the portal controller normalises the reference to a canonical \IRI\ that links to the industry platform's entry for that material. This design accommodates both public registries (ECHA, \PubChem) and industry platforms (ChemX, Materials Commons), treating them uniformly as external authoritative sources that \DPPs\ link to rather than duplicate. The Materials Commons initiative (Section~2.2) serves as one such external source. The portal-backed architecture complements Materials Commons by providing the storage, access governance, and ingestion validation that Materials Commons explicitly does not offer.

\subsection{The Open-Source Material Characterisation Database}

A conference participant described at \DPPforEU~2026~\cite{dpp4eu2026_day1} the development of ``\textit{an open-source database with the material characterisation for everybody to have access to identifiers for free and also to test material compositions}''. This initiative aligns with the cross-registry linking strategy: open-access material characterisation databases can serve as additional external authoritative sources that \DPPs\ link to. The portal controller's external entity normalisation can accommodate multiple registries --- ECHA for regulatory hazard data, \PubChem\ for chemical substance information, open-source material characterisation databases for materials property data --- normalising all references to canonical \IRIs\ and resolving them in real time during the Resolution Governance Loop. This multi-registry approach ensures that \DPPs\ carry the most authoritative, most current, and most comprehensive data available, without duplicating any of it.
